\chapter{Combined Torsional Analysis}
\label{ch:assignment5}

\begingroup\itshape Assignment 5 --- 17 points\endgroup

\bigskip

\section{Generalised torsion equation and boundary conditions}
\label{sec:a5-1}

\textbf{HET OPSTELLEN VAN HET PROBLEEM}

Uitwerking: differentiaalvergelijking, algemene oplossing per sectie, overgangsvoorwaarden bij de schotten en randvoorwaarden aan de uiteinden.

\finalanswer{\tbd{kort en expliciet antwoord}}

\section{Twist deformation and rate over the ship length}
\label{sec:a5-2}

\textbf{HET OPLOSSEN VAN HET PROBLEEM}

omdat phi' = 0 aan beide kanten van elk schot, ligt de vorm van de verdraaiing binnen elk blok al vast zonder dat je de buren nodig hebt. Het enige wat een blok doorgeeft aan het volgende is de hoek phi aan het eind.

Per blok ziet het er dan zo uit:

\begin{enumerate}
    \item Het moment legt C1 vast: C1 = $M_T$ / ($GI_p$).
    \item De twee voorwaarden aan de randen van het blok leggen C2 en C3 vast. Voor de middelste blokken is dat phi' = 0 links en rechts. Voor A en C is het phi'' = 0 aan het vrije uiteinde en phi' = 0 bij het schot.
    \item C0 volgt uit de hoek: bij blok A is dat phi(0) = 0, bij elk volgend blok de eindhoek van het vorige blok.
\end{enumerate}

\finalanswer{\tbd{kort en expliciet antwoord}}

\section{Open versus closed hatch covers}
\label{sec:a5-3}

\questionbox{Compare the results of a vessel subject to constant torsional
moment with all three hatch covers closed (BUH) and all open (BU) for
configuration B. \textbf{(5 pts)}
\begin{enumerate}[label=\alph*., leftmargin=*]
  \item Compare the plot of the rotation over the length, and discuss
        differences. \textbf{(2 pts)}
  \item Compare the plot of the twist rate over the length, and discuss
        differences. \textbf{(2 pts)}
  \item Compare the improvement with respect to the maximum twist rate.
        \textbf{(1 pt)}
\end{enumerate}}

\begin{enumerate}[label=\alph*., leftmargin=*]
  \item \tbd{Uitwerking}
  \item \tbd{Uitwerking}
  \item \tbd{Uitwerking}
\end{enumerate}

% \begin{figure}[htb]
%   \centering
%   \includegraphics[width=0.9\textwidth]{figures/<bestandsnaam>}
%   \caption{\tbd{onderschrift}}
%   \label{fig:a5BUBUH}
% \end{figure}

\finalanswer{\tbd{kort en expliciet antwoord per deelvraag}}
